\ficha{Mollo (1994), Controvérsias monetárias do século XIX}{mollo1994}

\begin{identificacao}
MOLLO, Maria de Lourdes Rollemberg. As controvérsias monetárias do século XIX.
\emph{Ensaios FEE}, Porto Alegre, v. 15, n. 1, p. 80-97, 1994.

Artigo de história do pensamento econômico, declaradamente um sumário para
iniciantes, 18 páginas, em português. Lido integralmente na versão publicada.
\end{identificacao}

\bloco{Problema e tese}
A autora sumaria as duas controvérsias monetárias inglesas porque a descrição
delas é dispersa e quase só existe em língua estrangeira (p. 80). A tese é
classificatória, não histórica. Um critério único ordena todas as posições, o
grau de inter-relação admitida entre a economia real e a monetária, e ele produz
gradação, não dicotomia (p. 93). Do lado ortodoxo, o privilégio da oferta de
moeda sustenta junto a exogeneidade e a neutralidade, e delas decorre a defesa
de regra fixa, o padrão-ouro entre elas. Do lado heterodoxo, a percepção da
demanda de moeda conduz à endogeneidade e à não neutralidade, e delas à recusa
da regra fixa (p. 93-94).

\bloco{Argumento}
\begin{enumerate}[leftmargin=*,nosep]
\item Primeira controvérsia, 1797 a 1825, travada sobretudo sob
inconversibilidade (p. 80). Bulionistas: Wheatley, Lauderdale, Ricardo e
Thornton. Antibulionistas: Bosanquet, Boase e Trotter (p. 81). Quatro questões
organizam a disputa e sobrevivem a ela: pertinência da Teoria Quantitativa da
Moeda e necessidade de controlar a oferta; relação entre a TQM e o balanço de
pagamentos; relação entre economia real e monetária; e relação entre moeda e
estabilidade de preços (p. 81).

\item O bulionista moderado já concede muito. Thornton aceita que o prêmio do
\emph{bullion} venha de drenagem interna, de restrição à exportação de ouro
amoedado ou da queda da libra dentro do custo de embarque, e aceita variação da
velocidade conforme o estado de confiança, o que torna o corte monetário em
crise potencialmente catastrófico (p. 81-82). É dele a crítica decisiva à
\emph{real bills doctrine}: as letras descontadas não dependem do produto físico
que as lastreia, mas da distância entre a taxa de juros bancária, travada em 5\%
pela lei da usura, e a taxa de lucro esperada (p. 83).

\item Ricardo é o polo máximo da ortodoxia por razão de método: analisa
equilíbrio de longo prazo, ignora o ajustamento de curto prazo, lê a taxa de
juros como fenômeno real e trata o crédito como mera transferência de renda
(p. 84). Wheatley e Lauderdale chegam a admitir efeitos reais ao discutir a
``poupança forçada'', mas o que os preocupa é a injustiça da redistribuição,
porque custos rígidos tornam a contração ``injuriosa'' para fazendeiros e
manufatureiros (p. 85). Os antibulionistas usam a mesma poupança forçada em
sentido oposto, como canal de renda nominal, gasto e emprego (p. 84-85).

\item Segunda controvérsia, 1825 a 1865. As duas escolas concordam no longo
prazo, valor da moeda pelo custo de produção do ouro e conversibilidade como
base da estabilidade (p. 86). A Currency School é mais restritiva que o
bulionismo: acrescenta à conversibilidade controles quantitativos de curto
prazo, como se a moeda fosse puramente metálica, daí o \emph{currency
principle}, a Palmer Rule de 1827 e a Lei Bancária de 1844 (p. 86-87). A Banking
School recusa cada peça, porque parte do ouro está entesourada e não circula,
porque a lei do refluxo dá limite natural à emissão, e porque depósitos e letras
também são meios de pagamento, de modo que o limite do departamento emissor é
compensado no departamento bancário (p. 88). Attwood e a Birmingham School vão
além e pregam o abandono da base metálica (p. 88).

\item O corte do item 2.2 é o que interessa reter. Para a Currency School a
moeda é véu: velocidade constante apaga a demanda de moeda, produto constante
decorre da neutralidade (p. 89). A Banking School, ao recuperar a demanda, chega
ao entesouramento e à determinação endógena da quantidade de moeda (p. 89), e as
suspensões da Lei Bancária em 1847, 1857 e 1866 parecem dar-lhe razão (p. 90).
Marx vai além dela: concorda com Tooke que o volume de moeda segue as
necessidades dos negócios, afirma a não neutralidade do crédito, que acelera a
produção e também as contradições que levam à crise, e trata a taxa de juros
como fenômeno monetário contra a taxa real de Ricardo e Overstone (p. 91-93).
\end{enumerate}

\chave{as preocupações quantitativas parecem reforçadas em função do acordo entre
as duas escolas sobre a necessidade de manter a conversibilidade-ouro das notas
bancárias, uma regra quantitativa fixa e automática de controle, que, na
controvérsia anterior, ainda dividia ortodoxos e heterodoxos}{p. 95}

\bloco{Conceitos e definições operacionais}
Neutralidade: a moeda não afeta variáveis reais, o que fixa o produto na equação
quantitativa (p. 89). Endogeneidade: a quantidade de moeda é determinada pelas
necessidades dos negócios e pelo entesouramento dos agentes, não pela decisão do
emissor (p. 89, p. 91). \emph{Currency principle}: a moeda em circulação deve
crescer quando entra ouro no país e cair quando sai (p. 87). Palmer Rule:
títulos em reserva constantes e emissão variando proporcional e diretamente com
o ouro em espécie (p. 87). Lei do refluxo: vencidos os empréstimos, as notas
voltam aos bancos e enxugam por si a circulação, o que dispensa controle
quantitativo (p. 88). \emph{Real bills doctrine}: o crédito bancário tem
contrapartida real porque nasce de letras ligadas a mercadorias em vias de venda
(p. 82-83). Poupança forçada: excesso de moeda que altera a distribuição em
favor do investimento, com preços de produtos subindo antes dos de fatores
(p. 84-85).

\bloco{Evidência e método}
Não há evidência empírica. É reconstrução de literatura secundária, com Viner
(1937) como espinha dorsal (p. 80, p. 82, p. 89) e apoio em Phyllis Deane e
Hayek. A única leitura de primeira mão é Marx, na edição brasileira de \emph{O
Capital} de 1971, com páginas dadas uma a uma (p. 91-93). Sem série, sem teste,
sem contrafactual, sem caso nacional fora da Inglaterra.

\bloco{Agentes e vertentes}
Bulionistas: Wheatley, Lauderdale, Ricardo e Thornton, este classificado como
moderado e mais heterodoxo que o grupo (p. 81). Antibulionistas: Bosanquet,
Boase e Trotter (p. 81). Currency School: Ricardo, Overstone e Torrens. Banking
School: Tooke e Fullarton (p. 86). Birmingham School: Attwood, o mais radical
(p. 88). A linhagem posterior é explicitada: Wicksell a partir de Thornton
(p. 83), Hayek a partir da poupança forçada (p. 85), Friedman e o banco central
independente como herdeiros da regra fixa (p. 93-94), Keynes e Minsky como
herdeiros da endogeneidade (p. 96).

\bloco{Críticas e limites}
Sumário didático que não abre disputa historiográfica, cita quase tudo de
segunda mão e arrola bibliografia que não reaparece no corpo. A gradação de
heterodoxia é critério do analista, não categoria dos debatedores, e mede
posições do início do século XIX por uma pergunta, a neutralidade, que só se
torna central depois. Há dois problemas de leitura: a passagem que invoca o
\emph{boom} de 1924-25 e a crise de 1926 para arbitrar entre as escolas (p. 90)
é anacrônica, e a virada deflacionária é datada de 1914 na conclusão (p. 94)
onde o argumento exige 1814, como está na página 80. Nada no artigo trata do
Brasil, de economia periférica ou de caixa de conversão, e nenhuma afirmação
sobre 1906 a 1914 pode ser apoiada nele.

\begin{dialogo}
\item Lei do refluxo no vocabulário de 1906. Procurar nos quatro jornais o
argumento de que a emissão da Caixa não pode ser excessiva porque as notas
voltam, porque são trocáveis por ouro ou porque a quantidade acompanha o volume
dos negócios. Mollo mostra que esse é o argumento da Banking School (p. 88).
Achá-lo na defesa da Caixa indica repertório heterodoxo dentro de instituição de
desenho ortodoxo.

\item Objeção dos meios de pagamento alternativos. Verificar se algum crítico
sustenta que o limite de emissão é ilusório porque o papel do Tesouro, os
depósitos bancários e o crédito comercial compensam o que a Caixa não emite. É a
objeção da Banking School à separação em dois departamentos (p. 88), e sua
presença testa se o debate discutia a eficácia do controle, e não só o nível da
taxa.

\item Registro da injustiça distributiva. Procurar, na crítica à valorização do
mil-réis, a fórmula de que a queda de preços redistribui arbitrariamente riqueza
e prejudica quem tem custos rígidos. Mollo situa esse registro nos bulionistas
moderados (p. 85), não na não neutralidade plena. Se for esse o argumento contra
a valorização, a posição é menos heterodoxa do que o rótulo papelista sugere.

\item Deslocamento do eixo com o consenso. Levantar se algum publicista de 1906
a 1914 ainda contesta a conversibilidade em si, ou se todos a aceitam e disputam
apenas taxa, limite e competência da Caixa. O paralelo com a Inglaterra depois
de 1825 (p. 95) dá a forma da hipótese: quando a regra vira consenso, a divisão
migra para os controles de curto prazo.
\end{dialogo}

\usonocapitulo{Parte III, como apoio à dobradiça, depois de Fonseca e Mollo
(2012). Serve a duas afirmações. Primeira, dá a definição precisa de cada
posição inglesa e das instituições que a Currency School propôs, o \emph{currency
principle}, a Palmer Rule e a Lei Bancária de 1844, o que permite ler o desenho
da Caixa de Conversão como parente institucional do \emph{currency principle}
sem atribuir a leitura à autora, que não trata do Brasil. Segunda, e mais
importante, dá o critério que separa duas perguntas embaralhadas pelo vocabulário
de época, a neutralidade da moeda e a endogeneidade da oferta, e mostra que elas
não andam sempre juntas, pois a Banking School insistia mais na impossibilidade
do controle do que na sua indesejabilidade (p. 95). A dobradiça está na
observação sobre a segunda controvérsia: quando a conversibilidade-ouro virou
acordo entre as escolas, deixou de dividir ortodoxos e heterodoxos, e a divisão
passou aos controles de curto prazo (p. 95). É o análogo estrutural do que o
capítulo sustenta para 1906, quando o eixo já não é metalismo contra papelismo,
mas valorização contra emissão e estabilidade à taxa nova.}
